\documentclass[12pt,a4paper]{article}
% =====================================================================
% Life Rays & Fragment Chains – Global Configuration (config.tex)
% =====================================================================

% Encoding, fonts, and language
\usepackage[utf8]{inputenc}
\usepackage[T1]{fontenc}
\usepackage[english]{babel}
\usepackage{lmodern}

% Page layout and spacing
\usepackage{geometry}
\geometry{a4paper, margin=2.5cm}

\usepackage{setspace}
\usepackage{parskip}

% Mathematics and graphics
\usepackage{amsmath}
\usepackage{amssymb}
\usepackage{amsfonts}
\usepackage{graphicx}
\graphicspath{{../figures/}}


% Section formatting
\usepackage{titlesec}
\titleformat{\section}{\normalfont\Large\bfseries}{\thesection}{1em}{}


% Header, footer, and page style
\usepackage{fancyhdr}
\usepackage{lastpage}


% Colors and hyperlink setup
\usepackage[dvipsnames]{xcolor}
\usepackage[colorlinks=true, linktoc=all, linkcolor=blue]{hyperref}

% Global hyperref configuration
\hypersetup{
    colorlinks=true,
    linkcolor=blue,
    filecolor=magenta,
    urlcolor=cyan,
    citecolor=ForestGreen,
    pdftitle={Life Rays \& Fragment Chains},
    pdfpagemode=FullScreen,
}

\usepackage[all]{hypcap}
\usepackage{tikz}
\usetikzlibrary{arrows.meta, positioning}


% Glossary system
\usepackage[nonumberlist]{glossaries}
\makeglossaries


% Index system
\usepackage{imakeidx}

% Custom index heading used in the Life Ray Theory style
\newcommand{\indexheading}[1]{%
\hypertarget{index:#1}{\centerline{\Large{#1}}\nopagebreak}%
\pdfbookmark[2]{#1}{index:#1}
}%
\makeindex[options=-s LifeRayTheory2.mst]


% --- Page Style Setup ---
\pagestyle{fancy}
\fancyhf{} % clear defaults

% --- Header ---
\fancyhead[L]{\textcolor{gray}{\small\textsc{Life Rays \& Fragment Chains}}}
\fancyhead[R]{\textcolor{gray}{\small\textit{Part II -- Formalisation  Fragments}}}

% --- Footer ---
\fancyfoot[L]{\textcolor{gray}{\footnotesize Independent Conceptual Research}}

% clickable DOI link:
\fancyfoot[R]{\footnotesize\href{https://doi.org/10.5281/zenodo.17561066}{\textcolor{gray!65!black}{DOI: 10.5281/zenodo.17561066}}}

% clickable page reference (x of y)
\fancyfoot[C]{\textcolor{gray}{\footnotesize Page \thepage\ of \pageref{LastPage}}}

% --- Line colors ---
\renewcommand{\headrulewidth}{0.4pt}
\renewcommand{\footrulewidth}{0.4pt}
\renewcommand{\headrule}{\hbox to\headwidth{\color{gray!60}\leaders\hrule height \headrulewidth\hfill}}
\renewcommand{\footrule}{\hbox to\headwidth{\color{gray!60}\leaders\hrule height \footrulewidth\hfill}}

\title{\textbf{Life Rays \& Fragment Chains --\\
A Theory of Encounter}}
\author{Peter Ngo}
\date{\today}

\begin{document}

\maketitle

\vspace{1.5em}
\begin{center}
\rule{0.45\textwidth}{0.4pt}\\[1em]

\textit{
For someone who once crossed my path,\\
whose presence became the origin of this theory.
}

\vspace{0.6em}

\textit{
And for the colleagues, friends, and companions\\
who helped transform intuition into structure\\
through challenge, questioning,\\
and refinement.
}

\rule{0.45\textwidth}{0.4pt}
\end{center}
\vspace{2em}

\tableofcontents
\newpage
\onehalfspacing

\section{Abstract}
This preprint develops the mathematical foundation of the Life Ray Theory, expanding the
conceptual framework introduced in Part~I into a formal system. The goal is to provide
the operator structures, notational conventions, and dynamical principles that allow
encounter, fragmentation, resonance, and coupling to be represented in a precise and
model-ready form. These formulations prepare the groundwork for later computational,
systemic, and empirical applications in Parts~III and~IV.


\vspace{1em}
\begin{center}
\textcolor{gray}{\small \textit{Preprint – Work in progress – Part I of a four-part research project.}}
\end{center}
\vspace{1em}

\section{Introduction}
Part~I of this project presented the phenomenological and conceptual architecture of
life rays, fragment chains, and shared spaces. Part~II marks a transition: it is the first
attempt to translate these intuitions into formal mathematical structures. The aim is not
to reduce the richness of encounter, but to render its mechanisms expressible through
operators, functions, and topological relations.

What follows is therefore a foundational step. It provides the notation, the structural
assumptions, and the core transformations upon which later models of resonance,
coupling, and dynamical evolution will rely. 

In this sense, the present text serves as the
bridge between the intuitive origins of the theory and its scientific articulation.

\subsection{Motivation and Scope}
Part~I introduced the core structures of the Life Ray Theory on a descriptive and
phenomenological level. Concepts such as life rays, fragments, shared space, and
resonant proximity were defined qualitatively and illustrated through structural
metaphors. Their function was conceptual clarity, not mathematical rigor.

The purpose of Part~II is fundamentally different. Its scope is to provide:
\begin{enumerate}
    \item a precise notational system for events, fragments, and spaces,
    \item formal operators for transduction, integration, and resonance,
    \item metrics for proximity, salience, and coupling,
    \item a formal salience function expressing weighted event imprinting,
    \item the topological assumptions that underlie the geometry of encounter,
    \item and a first formalization of resonant asymmetry (humor) as a
          special case of temporal divergence.
\end{enumerate}

The formulations introduced here are intentionally minimal: they aim to capture the
essential structure of interaction without binding the theory to a specific
computational model. Later parts will extend these foundations toward simulation,
multi-agent architectures, and empirical applications.

In short, the goal of this part is not to complete the theory, but to make it formally
expressible.

\section{Notation and Basic Objects}

Part~I introduced the intuitive architecture of life rays, fragments, and shared
spaces. In Part~II, these notions receive a formal representation. The following
definitions establish the basic mathematical objects on which all subsequent
operators and metrics will rely.

\subsection{Core Formal Structure}

We introduce the minimal formal structure underlying the Life Ray Theory.
All later operators, metrics, and interpretations are extensions of this core.

\subsection{Agents and Own Space}

Let $K$ be a set of agents.
Each agent $k \in K$ is associated with an own space $S_k$, representing its internal experiential state space.

\subsection{Life Rays}

A life ray is defined as a time-parametrized curve
\begin{equation}
\gamma_k : \mathbb{R}_{\ge 0} \to S_k,
\end{equation}
mapping personal time to internal state.

\subsection{Shared Events}

An encounter between two agents $i,j \in K$ at time $t$ is modeled as an event
\begin{equation}
E_{ij}(t).
\end{equation}
The shared space is defined as the set of all such events:
\[
C = \{ E_{ij}(t) \mid i,j \in K,\, t \in \mathbb{R}_{\ge 0} \}.
\]

\subsection{Fragments}

For each event and each involved agent, a fragment is generated by a transduction operator
\[
T_k : E_{ij}(t) \mapsto F_k.
\]
A fragment is minimally represented as
\begin{equation}
F_k = (t, v_k),
\end{equation}
where $t$ is the event time and $v_k$ an internal state imprint.

\subsection{Fragment Chains}

The fragment chain of agent $k$ at time $t$ is the ordered sequence
\begin{equation}
  \mathcal{F}_k(t) = (F_k^{(1)}, F_k^{(2)}, \dots, F_k^{(n)}),
  \label{eq:fragment-chain}
\end{equation}
containing all fragments generated up to time $t$.

This core defines the minimal ontology of the Life Ray Theory: agents, internal state spaces,
time-parametrized life rays, shared events, and fragment-based memory structures.
All further definitions refine this structure by adding semantic, affective, and dynamical detail.


\section{Extended Structures}

\subsection{Extended Event Structure}

While the core model represents an event as a minimal intersection $E_{ij}(t)$,
we now enrich events with descriptive and phenomenological parameters:

\begin{equation}
E_{ij} = (t, C, D, I, V, R),
\end{equation}

\subsection{Geometric Interpretation of Life Rays}

The life ray
\[
\gamma_k : \mathbb{R}_{\ge 0} \to \mathcal{S}_k
\]
may be interpreted both geometrically and temporally.

Geometrically, $\gamma_k$ represents a trajectory in the agent’s experiential state space
$\mathcal{S}_k$. Each point $\gamma_k(\tau)$ corresponds to the internal state of agent $k$
at personal time $\tau$. Intersections between $\gamma_i$ and $\gamma_j$ in shared space
correspond to encounter events $E_{ij}$.

Temporally, $\gamma_k$ is a parametrized curve whose parameter $\tau \in \mathbb{R}_{\ge 0}$
denotes the subjective time of the agent. In this sense, a life ray encodes the continuous
evolution of experience, while events represent discrete synchronization points between
different rays.

Two rays $\gamma_i$ and $\gamma_j$ may therefore intersect multiple times in shared space,
giving rise to a sequence of encounter events
\[
E_{ij}^{(1)}, E_{ij}^{(2)}, E_{ij}^{(3)}, \dots,
\]
as illustrated in Figure~\ref{fig:life_rays_events}.

\begin{figure}[h]
    \centering
        \begin{tikzpicture}[x=1cm,y=1cm,>=stealth,thick]
            % Time axis
            \draw[->,gray!70] (0,0) -- (10,0) node[right]{Time};

            % Life ray i
            \draw[blue] (0.5,1.2) -- (9.5,1.2)
                node[right,black]{\small Life ray $\gamma_i$};

            % Life ray j
            \draw[red] (0.5,-0.8) -- (9.5,-0.8)
                node[right,black]{\small Life ray $\gamma_j$};

            % Shared space window
            \fill[gray!15,rounded corners] (3.2,-1.3) rectangle (6.8,2.5);
            \node[gray!60] at (5.0,1.8) {\footnotesize shared space};

            % Intersection events
            \filldraw[blue!70] (3.8,1.2) circle (2.2pt);
            \filldraw[red!70]  (3.8,-0.8) circle (2.2pt);
            \draw[->,gray!70] (3.8,0.55) -- (3.8,0.95);
            \draw[->,gray!70] (3.8,-0.1) -- (3.8,-0.55);
            \node at (3.8,0.25) {\footnotesize $E_{ij}^{(1)}$};

            \filldraw[blue!70] (5.2,1.2) circle (2.2pt);
            \filldraw[red!70]  (5.2,-0.8) circle (2.2pt);
            \draw[->,gray!70] (5.2,0.55) -- (5.2,0.95);
            \draw[->,gray!70] (5.2,-0.1) -- (5.2,-0.55);
            \node at (5.2,0.25) {\footnotesize $E_{ij}^{(2)}$};

            \filldraw[blue!70] (6.2,1.2) circle (2.2pt);
            \filldraw[red!70]  (6.2,-0.8) circle (2.2pt);
            \draw[->,gray!70] (6.2,0.55) -- (6.2,0.95);
            \draw[->,gray!70] (6.2,-0.1) -- (6.2,-0.55);
            \node at (6.2,0.25) {\footnotesize $E_{ij}^{(3)}$};

        \end{tikzpicture}

    \caption{Two life rays ($\gamma_i,\gamma_j$) with repeated intersection events occurring within the shared space.}
    \label{fig:life_rays_events}
\end{figure}


\subsection{Fragments}

Every event $E_{ij}$ generates a perspective-dependent memory fragment
for each participant:
\[
F_k = (s_k, \phi_k, \theta_k, M_k),
\]
where $s_k$ is the salience, $\phi_k$ a semantic label,
$\theta_k$ the internal timestamp, and $M_k$ a set of affective, bodily,
or contextual markers.

\subsection{Fragment Chains}

The memory structure of an agent $k$ is modeled as a time-dependent
collection of discrete memory fragments.
Rather than assuming memory to be continuous, the Life Ray Theory
treats biographical memory as an ordered accumulation of such fragments.
At any time $t$, the fragment chain of agent $k$ is represented as an ordered sequence
defined in Eq.~\eqref{eq:fragment-chain}.

where each fragment corresponds to a previously transduced encounter event.
Together, these fragments define the evolving structure of memory over time.

We denote by $\mathbin{\|}$ the concatenation operator on fragment chains, such that
\begin{equation}
\mathcal{F}_k(t+\varepsilon) = \mathcal{F}_k(t) \mathbin{\|} F_k
\label{eq:fragment-concatenation}
\end{equation}
where $\mathbin{\|}$ appends the new fragment $F_k$ to the end of the existing ordered
sequence $\mathcal{F}_k(t)$.

Memory formation is assumed to be incremental rather than continuous.
When a new salient event occurs, it generates a fragment $F_k$ that is
added to the existing chain according to Eq.~\eqref{eq:fragment-concatenation}.
Here, $\varepsilon$ denotes a minimal temporal update step associated
with event processing, not a physical duration.

This formulation captures the idea that biographical memory grows
through the accumulation of discrete experiential units, while
preserving the temporal structure necessary for later integration,
resonance, and narrative coherence.

\subsection{Shared and Own Space}

We distinguish:
\begin{itemize}
    \item the \textbf{shared space} $C$, in which events $E_{ij}$ occur,
    \item and the \textbf{own space} $\mathcal{S}_k$, containing the
          internal representation and fragment chain of agent $k$.
\end{itemize}

An event $E_{ij}$ lies in shared space, while its transformed fragments
$F_i$ and $F_j$ reside in the respective own spaces.

\subsection{Time Notation}

We use:
\[
t \quad \text{(external event time)}, \qquad
\theta_k \quad \text{(internal time of fragment formation)},
\]
and, in the case of resonant asymmetry (see Section~\ref{sec:asymmetry}),
the distinction between expected and realized meaning times:
\[
t_e \neq t_r.
\]

These definitions provide the foundational vocabulary from which the
formal operators (transduction, integration, resonance), metrics
(salience, proximity, coupling), and geometric assumptions of later
sections can be expressed.

\section{Mechanism: Transduction}

\subsection{Definition}
An \emph{intersection event} $E_{ij}(t,C,D,I,V,R)$ arises in the
\emph{shared space} when the life rays of agents $i$ and $j$
interact at time $t$ within a context $C$ (location, task, social
setting), with duration $D$, intensity $I$, valence $V$, and
perceived reciprocity $R$.
\emph{Transduction} is defined as the operator
\begin{equation}
\mathcal{T}:\;
E_{ij}(t,C,D,I,V,R)
\;\mapsto\;
\big(F_i,\;F_j\big),
\end{equation}

\begin{figure}[h]
    \centering
        \begin{tikzpicture}[x=1cm,y=1cm,>=stealth,thick]

            % Event node
            \node[draw,rounded corners,fill=gray!10,align=center] (E) at (0,0)
                {\small Intersection event\\[-2pt] $E_{ij}(t,C,D,I,V,R)$};

            % Fragment nodes
            \node[draw,rounded corners,fill=blue!8,align=left] (Fi) at (6,1.0)
                {\small Fragment $F_i$:\\[-2pt]
                 \footnotesize content, timestamp, affect,\\[-3pt]
                 tags, label $\ell_i$, salience $s_i$};

            \node[draw,rounded corners,fill=red!8,align=left] (Fj) at (6,-1.0)
                {\small Fragment $F_j$:\\[-2pt]
                 \footnotesize content, timestamp, affect,\\[-3pt]
                 tags, label $\ell_j$, salience $s_j$};

            % Transduction arrows
            \draw[->] (E.east) -- (Fi.west)
                node[midway,above,sloped]{\footnotesize $\mathcal{T}$};
            \draw[->] (E.east) -- (Fj.west)
                node[midway,below,sloped]{\footnotesize $\mathcal{T}$};

            % Alignment (A)
            \draw[<->,gray!60,dashed] (Fi.south east) to[bend left=15]
                node[midway,right]{\footnotesize $\mathcal{A}(F_i,F_j)\in[0,1]$}
                (Fj.north east);

        \end{tikzpicture}

    \caption{Transduction: The shared event is transformed into two perspective-dependent fragments; an optional alignment $\mathcal{A}$ may occur through shared recall or co-narration.}
    \label{fig:transduction_diagram}
\end{figure}

which transforms the shared event into individual
\emph{fragments} $F_k$ (for $k\in\{i,j\}$).
A fragment minimally contains the fields
\[
F_k=\big(
\text{content},\,
\text{timestamp},\,
\text{affect},\,
\text{semantic tags},\,
\ell_k,\,
s_k
\big),
\]
where $\ell_k$ denotes a semantic label (short interpretation)
and $s_k\in\mathbb{R}_{\ge 0}$ represents the fragment’s
\emph{salience} (imprint strength).

\medskip
\subsection{Process Structure}
Transduction unfolds in three conceptual stages:
\begin{enumerate}
    \item \emph{Detection (Marking):}
    The event $E_{ij}$ exceeds an activation threshold $\theta$
    and is registered as salient.
    \item \emph{Encoding (Bundling):}
    Sensory input and evaluation are condensed into a
    meaning-bearing core, forming $F_k$
    (including tags and semantic label $\ell_k$).
    \item \emph{Integration (Chain Insertion):}
    The fragment $F_k$ is appended to the fragment chain of
    agent $k$; consolidation processes may increase or decrease
    $s_k$ over time.
\end{enumerate}

\medskip
\subsection{Asymmetry and Alignment}
Transduction is inherently perspective-dependent:
from the same event $E_{ij}$, two generally distinct fragments
$F_i$ and $F_j$ emerge, differing in semantic label $\ell$
and salience $s$.
Subsequent \emph{alignment} through shared recall or narration
can be modeled as a partial matching
\[
\mathcal{A}(F_i,F_j)\in[0,1],
\]
representing the degree of shared meaning.

\medskip
\subsection{Dynamics and Proximity (Helix)}
Repeated reciprocal intersection events reinforce existing
fragments or generate successor fragments.
At a macroscopic level, this process increases the
\emph{degree of intertwining} between life rays.
Operationally, a \emph{twist index} may be updated incrementally
after each event:
\[
\text{twist}_{ij}
\;\leftarrow\;
\text{twist}_{ij}
\;+\;
\beta_I\,\bar I
\;+\;
\beta_D\,D
\;+\;
\beta_R\,\bar R
\;+\;
\beta_F,
\]
where $\bar I$ and $\bar R$ aggregate both perspectives
(e.g.\ mean values), and $\beta_\bullet$ denote domain weights.

\subsection{Failure and Boundary Cases}
\begin{itemize}
    \item \emph{Sub-threshold events:}
    Insufficient intensity or novelty $\Rightarrow$ no stable fragment.
    \item \emph{Interference:}
    Competing events $\Rightarrow$ distorted or fragmented encoding.
    \item \emph{Reframing:}
    Later information alters contextual fit $\Rightarrow$
    re-weighting of $\ell_k$ and $s_k$.
\end{itemize}


\section{Salience and Fragment Integration}

\subsection{Salience Function}

The imprint strength of an event on an agent’s memory is quantified
by the \emph{salience} of the resulting fragment.
For an agent $k$, salience is modeled as a weighted sum:
\begin{equation}
s_k
= w_I I_k
+ w_D D
+ w_R R_k
+ w_F F_k^*
+ w_N N_k
- w_X X_k,
\quad w_\bullet \ge 0.
\end{equation}

Here, $I_k$ denotes the perceived interaction intensity,
$D$ the duration of the event,
$R_k$ the experienced reciprocity,
$F_k^*$ the frequency of similar prior events,
$N_k$ a novelty factor,
and $X_k$ represents interference or competing influences.
The weights $w_\bullet$ encode agent-specific sensitivity profiles.

This formulation reflects the assumption that not all encounters
contribute equally to memory formation.
Salience governs which fragments dominate the fragment chain,
thereby shaping memory depth, fragment density, and later resonance effects.

\subsection{Integration Operator}

While transduction generates individual fragments from shared events, biographical
coherence arises through their integration into the existing fragment chain.
We model this process by an integration operator acting on the memory structure of a
single agent $k$:

\begin{equation}
A_k : (\mathcal{F}_k(t), F_k) \longrightarrow \mathcal{F}_k(t+\varepsilon).
\end{equation}

In its minimal form, the operator reduces to sequence concatenation:
\begin{equation}
A_k(\mathcal{F}_k(t), F_k)
= \mathcal{F}_k(t) \mathbin{\|} F_k.
\end{equation}

The operator $A_k$ modifies the fragment chain by incorporating the new fragment
$F_k$ into the existing structure. Depending on contextual and semantic relations,
this may involve reinforcement, attenuation, or reweighting of existing fragments.


\section{Resonant Asymmetry (Humor)}
\label{sec:asymmetry}

\subsection{Phenomenological Origin}

In Part~I of the \emph{Life Ray Theory}, humor was introduced at a
phenomenological level.
Its central mechanisms were described conceptually, without formal
mathematical specification.

The core assumptions were the following:
\begin{itemize}
    \item Humor arises from an asynchrony between an expected meaning time
          $t_e$ and the actual realization time $t_r$.
    \item Prior to semantic resolution, the life ray exhibits a temporary
          deviation, perceived as a \emph{resonant asymmetry}.
    \item The delayed realization produces a local densification of the
          fragment structure, determining the subjective intensity of
          the humorous moment.
    \item The interpreting life ray conceptually \emph{winds} around the
          punchline’s origin before the effective intersection occurs.
\end{itemize}

This leads to the guiding intuition:

\begin{quote}
\textbf{Humor is a local increase of energetic density in the fragment
configuration, caused by temporally shifted intersection points of
two life rays.}
\end{quote}

In Part~I, this description was intentionally left without
formalization or visualization in order to preserve ontological clarity.

\subsection{Temporal Asymmetry}

The defining structural condition for humor is a temporal divergence
between expectation and realization:
\[
t_e \neq t_r.
\]
The greater this divergence, the more pronounced the pre-resolution
fragment accumulation becomes, up to a cognitive or affective saturation
threshold.

\subsection{Geometric Interpretation}

Formally, resonant asymmetry may be interpreted as a spiral-like
trajectory of the interpreting ray prior to intersection.
The number of windings reflects repeated prediction updates and
micro-fragment accumulation.
Let $W$ denote the number of effective windings before realization.

\subsection{First Humor Metric}

As a first approximation, the intensity of humor may be expressed as
\begin{equation}
A_{\text{humor}} \;\propto\; \frac{W}{\lvert t_r - t_e \rvert},
\end{equation}
where $W$ denotes the number of pre-intersection windings and
$\lvert t_r - t_e \rvert$ the temporal divergence between expectation and realization.

This expression captures two central intuitions:
(i) repeated interpretive cycling (larger $W$) amplifies the humorous effect, and
(ii) increasing temporal delay reduces coherence and thus diminishes humor.
In practice, the divergence of the expression for $\lvert t_r - t_e \rvert \to 0$ is
understood to be limited by cognitive saturation effects; consequently, this ratio
may be replaced in later models by a normalized or saturating function.

\subsection{Position within the Formal Framework}

With the introduction of salience and integration,
the Life Ray Theory becomes formally tractable:
events acquire operative structure,
fragments obtain a well-defined internal representation,
and interactions receive a mathematical signature.

Within this formal setting, humor constitutes a special dynamical case
of resonant interaction.
It emerges when high fragment salience coincides with controlled
temporal divergence.
In later parts, this formalization will support
the empirical parameterization of salience and temporal divergence
(Part~III) and the development of simulations and applied systems
(Part~IV).

\vspace{4em}
\noindent\rule{\textwidth}{0.4pt}
\vspace{1em}
\footnotesize
\noindent\textbf{Author’s Note and Disclosure} \\
This essay represents an independent conceptual work. All theoretical ideas, models,
definitions, and formal structures were conceived, developed, and decided exclusively
by the author.

An AI-based writing assistant (ChatGPT by OpenAI) was used solely for the purpose of
feedback and critical review, in particular with respect to clarity, structure,
consistency of notation, and linguistic formulation. The assistant was further employed
for minor language editing, translation, and formatting suggestions.

At no point did the AI system contribute original theoretical content, conceptual
design, or mathematical structure. All intellectual responsibility for the content,
interpretation, and validity of the presented theory remains entirely with the author.

\normalsize
\vspace{2em}

\newpage
\bibliographystyle{apalike}
\bibliography{references}

% --- Print Index ---
\clearpage
\phantomsection
\addcontentsline{toc}{section}{Index}

\printindex
\end{document}
